\section{引言}

糖尿病是一类以慢性高血糖为主要表现的代谢性疾病。
长期血糖控制不佳会增加心血管、肾脏、神经和视网膜等并发症风险。
世界卫生组织指出，糖尿病仍是重要公共卫生问题；
Reza 等的研究也强调，早期诊断和管理有助于预防或延缓并发症发生
\cite{whoDiabetesFactSheet2024,rezaImprovingDiabetesDisease2024}。
在常规筛查场景中，结构化体检指标比复杂医学检查更容易获取，
因此利用机器学习模型进行糖尿病风险预测具有实际意义。

本文使用 Pima Indians Diabetes Database 开展二分类实验。
该数据集包含 768 名样本、8 个数值型输入特征和 1 个二分类标签，
样本均为至少 21 岁的 Pima Indian 女性\cite{uci_pima}。
该数据集规模较小但变量清晰，
适合在统一预处理流程下比较不同分类算法的性能和误判特征。

本文围绕糖尿病风险预测流程展开：
数据层面完成缺失值、异常值、标准化和分层划分处理；
模型层面以 Logistic Regression 作为可解释基准，
并在相同实验设置下比较 KNN、Decision Tree、Random Forest、SVM、Gradient Boosting、XGBoost 和 LightGBM 等多种模型；
结果分析部分结合分类指标、ROC AUC、混淆矩阵和特征解释评估模型表现，
并将最终模型接入交互式演示系统。
